% Source: https://artofproblemsolving.com/community/c3h2325113_combinatorics_is_hard
% Stable identity: aops-topic-2325113-q01
Read the shelves in order, first shelf left to right, then the second, then the third. This turns an arrangement into a sequence of $12$ colours together with two ``cuts'' that split it into three nonempty consecutive parts. Conversely, each sequence with two cuts gives exactly one arrangement, so this is a bijection.

\begin{itemize}
\item The number of colour sequences is the multinomial coefficient
\[\dfrac{12!}{5!\,3!\,3!\,1!}=110880.\]
\item The two cuts go into two of the $11$ gaps between consecutive objects (different gaps, so that no shelf is empty): $\dbinom{11}{2}=55$ ways.
\end{itemize}
The total is $110880\cdot55=\boxed{6098400}$.
